\documentclass[../../../include/open-logic-section]{subfiles}

\begin{document}

\olfileid{sth}{choice}{woproblem}
\olsection{مشكلة الترتيب الجيد}

لا يخفى أن قبول الترتيب الجيد تتعلق به نتائج كثيرة. غير أن البحث في
هذا الأصل قد جرى غالبًا على أصل يكافئه، وهو مسلّمة الاختيار. فلنصرف
النظر إليها الآن، ثم نقيم البرهان على التكافؤ.

في~\citeyear{Cantor1883} أظهر كانتور نصرته لمسلّمة الترتيب الجيد،
ووصفها بأنها «قانون للفكر يبدو لي أساسيًا، غنيًا بنتائجه، وجديرًا
بملاحظة خاصة لعمومية صلاحيته»؛ وهذا بنقل~\citeauthor{Potter2004}
\citeyear[p.~243]{Potter2004}. ثم انتهى كانتور إلى أن هذه «المسلّمة»
نفسها محتاجة إلى برهان، ووافقه جمهور الرياضيين.

وزعم زيرميلو أنه حل الإشكال في~\citeyear{Zermelo1904}. ولا بد لفهم
حله من تعريفات.

\begin{defn}
تكون الدالة~${\OLId{latin-ordinary}{f}}$ \emph{دالة اختيار} إذا وفقط إذا كان~${\OLId{latin-ordinary}{f}}({\OLId{latin-ordinary}{x}}) \in {\OLId{latin-ordinary}{x}}$
لكل~${\OLId{latin-ordinary}{x}} \in \dom{{\OLId{latin-ordinary}{f}}}$. وتكون~${\OLId{latin-ordinary}{f}}$ \emph{دالة اختيار لـ${\OLId{latin-looped}{A}}$} إذا وفقط
إذا كانت~${\OLId{latin-ordinary}{f}}$ دالة اختيار و$\dom{{\OLId{latin-ordinary}{f}}} = {\OLId{latin-looped}{A}} \setminus \{\emptyset\}$.
\end{defn}

ومعنى ذلك في التصور أن دالة الاختيار لـ~${\OLId{latin-looped}{A}}$ تعين، من كل مجموعة غير
خالية~${\OLId{latin-ordinary}{x}} \in {\OLId{latin-looped}{A}}$، عنصرًا مخصوصًا~${\OLId{latin-ordinary}{f}}({\OLId{latin-ordinary}{x}})$ ينتمي إلى~${\OLId{latin-ordinary}{x}}$. فمسلّمة
الاختيار هي:

\begin{axiom}[الاختيار]
	لكل مجموعة دالة اختيار.
\end{axiom}

وقد بين زيرميلو أن الاختيار يستلزم الترتيب الجيد، وأن العكس كذلك.

\begin{thm}[في $\ZF$]\ollabel{thmwochoice}
الترتيب الجيد والاختيار متكافئان.
\end{thm}

\begin{proof}
\emph{من اليسار إلى اليمين.} لتكن~${\OLId{latin-looped}{A}}$ مجموعة من المجموعات. توجد
$\bigcup {\OLId{latin-looped}{A}}$ بمسلّمة الاتحاد؛ وبالترتيب الجيد يوجد ترتيب~$<$ يجعل
$\bigcup {\OLId{latin-looped}{A}}$ مرتبة ترتيبًا جيدًا. ثم ضع
${\OLId{latin-ordinary}{f}}({\OLId{latin-ordinary}{x}}) = \text{أصغر عنصر في ${\OLId{latin-ordinary}{x}}$ بحسب $<$}$. فتكون هذه دالة اختيار
لـ~${\OLId{latin-looped}{A}}$.

\emph{من اليمين إلى اليسار.} خذ~${\OLId{latin-looped}{A}}$. إن كانت~${\OLId{latin-looped}{A}} = \emptyset$، كان
الترتيب الخالي ترتيبًا جيدًا لها وتم الأمر. فلتكن إذن
${\OLId{latin-looped}{A}} \neq \emptyset$. يقتضي الاختيار وجود دالة اختيار~${\OLId{latin-ordinary}{f}}$
لـ~$\Pow{{\OLId{latin-looped}{A}}} \setminus \{\emptyset\}$. ثم عرّف بالعود ما فوق المتناهي
حدًا على الوجه الآتي:
\begin{align*}
	{\OLId{latin-ordinary}{g}}({\OLMathNumeral{0}}) &= {\OLId{latin-ordinary}{f}}({\OLId{latin-looped}{A}})\\
	{\OLId{latin-ordinary}{g}}({\OLId{greek-variable}{alpha}}) &=
		\begin{cases}
			\text{توقف!{}} &\text{إذا كان }{\OLId{latin-looped}{A}} = \funimage{{\OLId{latin-ordinary}{g}}}{{\OLId{greek-variable}{alpha}}}\\
			{\OLId{latin-ordinary}{f}}({\OLId{latin-looped}{A}} \setminus \funimage{{\OLId{latin-ordinary}{g}}}{{\OLId{greek-variable}{alpha}}}) & \text{خلاف ذلك}\\	
		\end{cases}
\end{align*}
ولفظ «توقف!» رمز مختصر عن تعريف يطول بسطه: إذا تحقق
${\OLId{latin-looped}{A}} = \funimage{{\OLId{latin-ordinary}{g}}}{{\OLId{greek-variable}{alpha}}}$ أول مرة، جعلنا
${\OLId{latin-ordinary}{g}}(\delta) = {\OLId{latin-looped}{A}}$ لكل~$\delta \geq {\OLId{greek-variable}{alpha}}$. ولا يكفل لنا أول الأمر إلا
تعريف حد، كما تبين الملاحظات بعد
\olref[sth][spine][recursion]{transrecursionschema}؛ وسنثبت أنه دالة.

لما كانت~${\OLId{latin-ordinary}{f}}$ دالة اختيار، فلكل~${\OLId{greek-variable}{alpha}}$ يسبق أول موضع للتوقف يكون
${\OLId{latin-ordinary}{g}}({\OLId{greek-variable}{alpha}}) = {\OLId{latin-ordinary}{f}}({\OLId{latin-looped}{A}} \setminus \funimage{{\OLId{latin-ordinary}{g}}}{{\OLId{greek-variable}{alpha}}}) \in {\OLId{latin-looped}{A}} \setminus
\funimage{{\OLId{latin-ordinary}{g}}}{{\OLId{greek-variable}{alpha}}}$؛ أي~${\OLId{latin-ordinary}{g}}({\OLId{greek-variable}{alpha}}) \notin \funimage{{\OLId{latin-ordinary}{g}}}{{\OLId{greek-variable}{alpha}}}$.
فإذا سبق~${\OLId{greek-variable}{alpha}}$ و$\beta$ أول موضع للتوقف، وكان
${\OLId{latin-ordinary}{g}}({\OLId{greek-variable}{alpha}}) = {\OLId{latin-ordinary}{g}}(\beta)$، لزم~${\OLId{latin-ordinary}{g}}(\beta) \notin
\funimage{{\OLId{latin-ordinary}{g}}}{{\OLId{greek-variable}{alpha}}}$، أي~$\beta \notin {\OLId{greek-variable}{alpha}}$، وكذلك
${\OLId{greek-variable}{alpha}} \notin \beta$. ومن التثليث~${\OLId{greek-variable}{alpha}} = \beta$. فقيد~${\OLId{latin-ordinary}{g}}$ على
القيم السابقة للتوقف !!{injective}.

ثم لا بد أن يقع التوقف فعلًا؛ أي يوجد أصغر عدد ترتيبي~${\OLId{greek-variable}{alpha}}$ بحيث
${\OLId{latin-looped}{A}} = {\OLId{latin-ordinary}{g}}[{\OLId{greek-variable}{alpha}}]$. فإن لم يكن، كانت~${\OLId{latin-ordinary}{g}}$ !!{injective} على كل عدد
ترتيبي، فحصلنا على~$\cardless{{\OLId{greek-variable}{alpha}}}{{\OLId{latin-looped}{A}}}$ لكل عدد ترتيبي~${\OLId{greek-variable}{alpha}}$،
وذلك يناقض~\olref[hartogs]{HartogsLemma} مطبقة على~${\OLId{latin-looped}{A}}$. وعند أول
مرحلة توقف~${\OLId{greek-variable}{alpha}}$، يكون
${\OLId{latin-ordinary}{g}}\restrict{\OLId{greek-variable}{alpha}} \colon {\OLId{greek-variable}{alpha}} \to {\OLId{latin-looped}{A}}$ !!a{bijection}؛ لأن قيمه قبل
التوقف متمايزة و${\OLId{latin-ordinary}{g}}[{\OLId{greek-variable}{alpha}}] = {\OLId{latin-looped}{A}}$.

فمن مجموع ذلك يستعمل~${\OLId{latin-ordinary}{g}}\restrict{\OLId{greek-variable}{alpha}}$ في ترتيب~${\OLId{latin-looped}{A}}$ ترتيبًا جيدًا.
\end{proof}

فالترتيب الجيد والاختيار لا يثبت أحدهما دون صاحبه ولا يسقط دونه. ويبقى
بعد هذا: أيثبتان أم يسقطان؟

\end{document}
